% !TEX root = /home/eb/.vscode/extensions/qftrules.programme-colle/tmp/Exercice.tex
\graphicspath{{\pathtoexos/Info/Figures/}}

%%%%%%%%%%%%%%%%%%%%%%%%%%%%%%%%%%%
\begin{exo}[PC][2][devoir]{Classification d'iris}
\Enonce{
  Les fleurs iris sont des plantes à fleurs très appréciées pour leur beauté et leur diversité. Elles appartiennent à la famille des iridacées et sont caractérisées par leurs pétales colorés et leurs formes élégantes. 
Les trois variétés d'iris les plus connues sont l'iris setosa, l'iris versicolor et l'iris virginica. Ces trois variétés présentent des différences subtiles dans la forme et la taille de leurs pétales et de leurs sépales, comme illustré sur la \Fig{iris}. Des mesures de la largeur $\ell$ et de la longueur $L$ des pétales ont été effectuées sur un ensemble de $N=150$ iris, représentée par une liste de points du plan appelée \texttt{iris}. Ces mesures sont représentées dans le plan par les points $(\ell,L)$ sur la \Fig{iris_graphe}.
\begin{fig}[]
  \centering
  \begin{subfig}[iris][0.5]
    \includegraphics[height=0.53\columnwidth]{iris.png}
  \end{subfig}\hfill
  \begin{subfig}[iris_graphe][0.45]
    \includegraphics[height=0.68\columnwidth]{iris_graphe.png}
  \end{subfig}
  \capt{(a) Trois variétés d'iris : setosa, versicolor et virginica.
  (b) Représentation des données iris dans le plan $(\ell,L)$, où $\ell$ est la largeur et $L$ la longueur des pétales.
  }
\end{fig}
}

\begin{quest}
  \QR	{Rappeler la différence entre un apprentissage supervisé et un apprentissage non supervisé. Quel algorithme de classement standard est associé à chaque type d'apprentissage ?}
  {
  Voir cours : apprentissage supervisé = données étiquetées (algorithme des $k$ plus proches voisins),
  apprentissage non supervisé = données non étiquetées (algorithme des $k$-moyennes).
  }
\end{quest}

\subsection{Algorithme des k plus proches voisins}

\begin{quest}
  \QR	{Définir mathématiquement la distance euclidienne $d(y_1,y_2)$ entre deux points $y_1=(\ell_1,L_1)$ et $y_2=(\ell_2,L_2)$ du plan.  La définir en \texttt{python} en complétant le code ci-dessous et en utilisant la fonction \texttt{np.sqrt} de la bibliothèque \texttt{numpy}. Quelle autre distance connaissez-vous ?
  \begin{minted}[linenos=true,frame=single]{python}
    def distance(y1,y2):
      """
      Calcule la distance euclidienne entre deux paires de coordonnées y1 et y2.
      """
      return ...
  \end{minted}
  }
  {
    Mathématiquement,
    \eqbox{
      d(y_1,y_2) = \sqrt{(\ell_1-\ell_2)^2 + (L_1-L_2)^2}.
    }
    En python,
  \inputminted[linenos=true,frame=single]{python}{\pathtoexos/Info/Scripts/ia_distance_euclid.py}
    }

  \QR {On dispose de la classification \texttt{varietes} (de type dictionnaire) qui associe à chaque iris sa variété. Les classes correspondent aux chaînes de caractère \texttt{`setosa'}, \texttt{`versicolor'} et \texttt{`virginica'}. Que renvoie \texttt{varietes[(6,1.8)]} en supposant que \texttt{(6,1.8)} est bien un élément de la liste \texttt{iris} ?}
  {
L'appel \texttt{varietes[(6,1.8)]} renvoie la variété \boxed{\texttt{`virginica'}}.
  }
  
  \QR {Selon vous, où devrait être classé un nouvel iris de coordonnées $(\ell=2.5,L=0.5)$ ? de coordonnées $(\ell=3.5,L=1.7)$ ? Justifier. }
{
Le point $(2.5,0.5)$ est proche des iris de la variété setosa, tandis que le point $(3.5,1.7)$ est proche des iris de la variété versicolor. Il est donc probable que le premier iris soit classé dans la variété \boxed{\texttt{'setosa'}} et le second dans la variété \boxed{\texttt{'versicolor'}}.
}

  \QR {On souhaite appliquer l'algorithme des $k$ plus proches voisins à l'ensemble d'apprentissage \texttt{iris} et à la classification \texttt{varietes} pour classer un nouvel iris de coordonnées $(\ell=2.5,L=0.5)$. On commence pour cela par définir la fonction \texttt{voisins} suivante. Qu'effectue cette fonction ? Effectuer l'appel pas à pas de cette fonction sur l'ensemble \texttt{E = [(0,0),(1,1),(2,2),(3,3),(4,4)]}, \texttt{x=(1.9,1.9)}, \texttt{k=3} et \texttt{d=distance}.
  \begin{minted}[linenos=true,frame=single]{python}
def voisins(x, E, k, d):
    indices = sorted(E, key=lambda y: d(x, y))
    return indices[:k]
  \end{minted}}
  {
    La fonction \texttt{voisins} renvoie les $k$ points les plus proches voisins de la donnée $x$ dans l'ensemble $E$. La liste indices contient les points de $E$ triés par ordre croissant de leur distance à $x$.
  }

    \QR {On définit maintenant la fonction \texttt{classemaj} que renvoie la classe majoritaire de la liste des $k$ plus proches voisins d'une nouvelle donnée $x$. Préciser les spécifications de cette fonction (variables d'entrée, variables de sortie). Effectuer l'exécution pas à pas de l'appel de cette fonction avec les arguments \texttt{E = [(0,0),(1,1),(2,2),(3,3),(4,4)]}, \texttt{x=(1.9,1.9)}, \texttt{k=3}, \texttt{d=distance}, et \texttt{p=\{(0,0):`a',(1,1):`b',(2,2):`b',(3,3):`a', (4,4):`b'\}}. Quelle classe majoritaire est prédite sur cet exemple simple ?
  \begin{minted}[linenos=true,frame=single]{python}
  def classemaj(E, x, d, k, p):
    v = voisins(E, x, d, k)             # k plus proches voisins de x
    c = {p[y]:0 for y in v}             # liste des classes des voisins
    for y in v:                         # boucle sur voisins
        c[p[y]] += 1                    # comptage des classes
    maxi = 0                            # initialisation maximum
    for i in c:                         # boucle sur classes des voisins
        if c[i] > maxi:                 # si classe i plus nombreuse
            maxi, classe = c[i], i      # mise à jour classe majoritaire
    return classe                       # retourne classe majoritaire
  \end{minted}}
  {
    % \begin{minted}[linenos=true,frame=single]{python}
    %   def classemaj(E, x, d, k, p):
    %     """
    %     Entrée : E (list) : ensemble d'apprentissage
    %              x (tuple) : nouvelle donnée
    %              d (func) : distance
    %              k (int) : nombre de voisins
    %              p (dict) : classification
    %     Renvoie la classe majoritaire des k plus proches voisins de x dans E.
    %     Sortie : classe (str) : classe majoritaire
    %     """
    %     ...
    %   \end{minted}
      L'appel de la fonction \texttt{classemaj} avec les arguments donnés renvoie la classe majoritaire \boxed{\texttt{`b'}}.
  }

  \QR {En appliquant la fonction \texttt{classemaj} à l'ensemble d'apprentissage \texttt{E=iris} et à la classification \texttt{p=varietes}, sur la nouvelle donnée $(\ell=2.5,L=0.8)$, on obtient les résultats suivants suivant la valeur de $k$ : `setosa' pour $k=1$, `versicolor' pour $k=3$, `setosa' pour $k=5$, `setosa' pour $k=7$. Commenter ces résultats.  Pourquoi seules les valeurs de $k$ impaires sont-elles utilisées ?}
  {
    Pour $k=1$, la classe majoritaire est `setosa' car le point $(2.5,0.8)$ est proche d'un seul iris de cette variété. Pour $k=3$, la classe majoritaire est `versicolor' car le point $(2.5,0.8)$ est proche de deux iris de cette variété et d'un iris de la variété `setosa'. Pour $k=5$, la classe majoritaire est `setosa' car le point $(2.5,0.8)$ est proche de trois iris de cette variété et de deux iris de la variété `versicolor'. Pour $k=7$, la classe majoritaire est `setosa' car le point $(2.5,0.8)$ est proche de quatre iris de cette variété et de trois iris de la variété `versicolor'. La classe prédite dépend du nombre de voisins ! Les valeurs de $k$ impaires sont utilisées pour éviter les cas d'\fbox{égalité de voisinage}.
  
  }
\end{quest}

\subsection{Qualité des prédictions}

\Enonce{
Les données de l'ensemble \texttt{iris} précédent sont déjà classéEs dans le dictionnaire \texttt{varietes} et peuvent servir d'ensemble test. On souhaite quantifier la qualité du classement par l'algorithme des $k$ plus proches voisins. On définit alors la matrice de confusion $M$ associée à cet ensemble test, telle que $M_{ij}$ est le nombre d'iris de la classe réelle $i$ dont la classe attribuée par l'algorithme des $k$ plus proches voisins est $j$. Sur les $N = 150$ iris, pour une certaine valeur de $k$, on obtient la matrice de confusion de la \Tab{confusion}.
% \begin{center}
  \begin{table}[H]
    \centering
    \begin{tabular}{|c|c|c|c|}
      \hline
      & setosa & versicolor & virginica \\
      \hline
      setosa & 46 & 3 & 0 \\
      versicolor & 1 & 48 & 2 \\
      virginica & 0 & 5 & 45 \\
      \hline
    \end{tabular}
    \caption{Matrice de confusion.}
    \label{tab:confusion}
    
  \end{table}
% \end{center}
}

  \begin{quest}
    \QR	{La précision $P$ de l'algorithme des $k$ plus proches voisins est définie comme le nombre d'iris bien classés sur le nombre total d'iris. Calculer la précision $P$ sur la matrice de confusion précédente.}
    {
      La précision $P$ est donnée par le rapport du nombre d'iris bien classés sur le nombre total d'iris, soit
      \eqbox{
        P = \frac{46+48+45}{150} = \frac{139}{150} \approx 0.93.
      }
    }

    \QR {Implémenter une fonction \texttt{precision} qui prend comme argument une matrice de confusion $M$ de taille $3 \times 3$ et qui renvoie la précision $P$. On précisera les spécifications de cette fonction et on la commentera. Exécuter cette fonction pas à pas sur la matrice de confusion de la \Tab{confusion}.}
    {
      Voici une version rapide. Voir cours pour version plus longue (et plus simple).
        \inputminted[linenos=true,frame=single]{python}{\pathtoexos/Info/Scripts/ia_precision_rapide.py}
      La fonction \texttt{precision} renvoie la précision de l'algorithme des $k$ plus proches voisins. Pour la matrice de confusion donnée, on a \texttt{precision(M) = 0.93}.
    }

    \QR {On représente sur la \Fig{precision} le graphe de la précision $P(k)$ en fonction du nombre de plus proches voisins $k$. Commenter l'allure de cette courbe. Quelle valeur de $k$ est optimale pour ce jeu de données ?}
    {
      La précision $P$ diminue avec le nombre de voisins $k$ pour les valeurs de $k$ suffisamment grandes,
      car le nombre de voisins devient trop important et l'algorithme perd en spécificité. La précision augmente avec $k$ pour les valeurs de $k$ suffisamment petites, car l'algorithme gagne en robustesse. La précision est maximale pour une valeur intermédiaire de $k$. La valeur optimale de $k$ est celle qui maximise la précision, soit \boxed{$k=10$}.
    }
    % \item On peut également utiliser le taux d'erreur pour évaluer la qualité de l'algorithme des $k$ plus proches voisins. Le taux d'erreur $TE$ est défini comme le nombre d'iris mal classés sur le nombre total d'iris. Calculer le taux d'erreur $TE$ sur la matrice de confusion précédente.

      \QR	{Pour sonder la qualité de l'algorithme des $k$ plus proches voisins, on souhaite visualiser la frontière de décision. 
      
      La \emph{frontière de décision} d'une classe $i$ est définie comme la courbe $L_{i}(\ell)$ qui sépare deux régions du plan $\mathbb{R}^2$ : la région $R_i$ telle que pour tout point $y \in R_i$, la classe prédite est $i$, et la région $\overline{R_i}$ telle que pour tout point $y \in \overline{R_i}$, la classe prédite n'est pas $i$.
      
      On visualise les frontières de décision pour $k=3$ sur la \Fig{decision}. Quelle classe serait prédite pour le point de coordonnées $(\ell=3,L=3)$ ?}
      {
        Le point $(3,3)$ est dans la région de décision de la classe versicolor. Il serait donc classé dans la variété \boxed{\texttt{`versicolor'}}. Sans tracer les frontières de décision, il serait moins évident de classer ce point à l'œil nu !
      }
    \end{quest}

    \Enonce{
    \begin{fig}[label]
      \centering
      \begin{subfig}[precision][0.48]
        \includegraphics[width=1.15\columnwidth]{precision_k.pdf}
      \end{subfig}
      \begin{subfig}[decision][0.48]
        \includegraphics[width=1.15\columnwidth]{decision_boundary.pdf}
      \end{subfig}
      \capt{
        (a) Précision $P$ de l'algorithme des $k$ plus proches voisins en fonction du nombre de voisins $k$.
        (b) Frontières de décision entres les trois régions bleu (setosa), rouge (versicolor) et verte (virginica) pour $k=3$.
      }
    \end{fig}
    }

\subsection{Algorithme des $k$-moyennes}

\Enonce{
On suppose maintenant que les données de l'ensemble \texttt{iris} ne sont pas étiquetées, mais on sait que chaque fleur appartient à l'une des trois variétés d'iris : setosa, versicolor et virginica. On souhaite donc les classer par l'algorithme des $k$-moyennes avec $k=3$. On rappelle le principe de l'algorithme des $k$-moyennes dans l'encadré ci-dessous.
\begin{center}
  
  \fbox{%
  \parbox{\dimexpr\linewidth-2\fboxsep-2\fboxrule\relax}{
    \begin{enumerate}[label=(\Alph*)]
      \item On choisit initialement $k$ centres de classes $c_1$, ..., $c_k$.
      \item On associe à chaque donnée $x$ la classe $i$ telle que $d(x,c_i)$ est minimal.
      \item On évalue les barycentres $b_1$, ..., $b_k$ des classes, que l'on choisit comme nouveaux centres.
    \end{enumerate}
    }}
  \end{center}
\udl{Étape (A)}. Il existe plusieurs façons d'initialiser les centres : centres aléatoires dans l'espace des paramètres, centres aléatoires parmi les données, ou centres déterminés par une méthode heuristique. Un exemple de choix heuristique est le suivant : le premier centre $c_1$ est choisi aléatoirement parmi les données ; chaque centre suivant $c_i$ est choisi comme la donnée la plus éloignée des centres $c_1$, ..., $c_{i-1}$ déjà définis.
}

\begin{quest}
  \QR	{On définit ci-dessous une fonction \texttt{eloigne} qui renvoie la donnée $x$ dont la somme des distances $d(x,c)$ aux centres $c$ est maximale. Préciser les spécifications de cette fonction et expliquer la syntaxe de ce code.
  % \sol{
  \begin{minted}[linenos=true,frame=single]{python}
    def eloigne(liste, E, d):
        return max(E, key=lambda x: sum(d(x,c) for c in liste))    
  \end{minted}
  % }
  }
  {
    La fonction \texttt{eloigne} renvoie la donnée $x$ de l'ensemble $E$ dont la somme des distances aux éléments de la liste est maximale. La fonction \texttt{max} renvoie l'élément de $E$ qui maximise la fonction \texttt{lambda x: sum(d(x,c) for c in liste)}. La fonction \texttt{sum}  renvoie la somme des distances de $x$ aux centres de la liste. La méthode \texttt{lambda} permet de définir une fonction anonyme.
  }

  \QR {Écrire une fonction \texttt{init\_centres} qui prend comme arguments un entier $k$, un ensemble de données \texttt{E}, une distance \texttt{d}, et qui renvoie une liste de $k$ centres initialisés selon la méthode heuristique décrite ci-dessus. On précisera les spécifications de cette fonction et on la commentera. Quelles seraient approximativement les positions des centres $c_2$ et $c_3$ si le premier centre est choisi aléatoirement en $c_1=(1.5,0.5)$ ? Reproduire l'allure du graphe de la \Fig{iris_graphe} en ajoutant de façon visible les points $c_1$, $c_2$ et $c_3$ pour illustrer votre réponse.}
  {
      \inputminted[linenos=true,frame=single]{python}{\pathtoexos/Info/Scripts/ia_init_centres.py}
    Les centres $c_2$ et $c_3$ seraient approximativement en $(4.5,1.5)$ et $(5.5,2.5)$.
  
  }
\end{quest}

 
\begin{quest}
  \QR	{ \udl{Étape (B)}. On définit la fonction \texttt{minclasse} qui prend comme arguments une donnée \texttt{x}, une liste de centres \texttt{centres}, et une distance \texttt{d}, et qui renvoie le centre $c$ le plus proche de \texttt{x}. Compléter la boucle \texttt{for} du code suivant pour obtenir cette fonction.
  % \begin{minted}[linenos=true,frame=single]{python}
  % def minclasse(x, centres, d):
  %   min(centres, key = lambda y:d(x,y))
  %   \end{minted}
\begin{minted}[linenos=true,frame=single]{python}
def minclasse(x, centres, d):
  """
  Entrée : x (tuple) : une donnée
            centres (list) : une liste de centres
            d (func) : une distance
  Renvoie le centre le plus proche de x parmi les centres.
  Sortie : c (tuple) : le centre le plus proche de x
  """
  c    = centres[0]
  dist = d(x,c)
  for y in centres[1:]:
      ...
  return c
\end{minted}}
  {
    \inputminted[linenos=true,frame=single]{python}{\pathtoexos/Info/Scripts/ia_minclasse_sol.py}
  }

  \QR {On définit ensuite la fonction \texttt{partition} de la façon suivante.  Que réalise cette fonction ? Commenter le code et expliquer chaque étape.
    \inputminted[linenos=true,frame=single]{python}{\pathtoexos/Info/Scripts/ia_partition_minclasse.py}}
  {
  
  }
\end{quest}


\begin{quest}
  \QR	{\udl{Étape (C)}.  Rappeler la définition du barycentre $b$ d'un ensemble de points $x_i \in \mathbb{R}^2$, pour $i=1,...,n$. Implémenter une fonction \texttt{barycentre} qui prend comme argument une liste de points du plan, appelée \texttt{points}, et qui renvoie le barycentre de ces points. On précisera les spécifications de cette fonction et on la commentera.}{

  }
\end{quest}

\Enonce{
On souhaite maintenant appliquer l'algorithme des $k$-moyennes à l'ensemble de données \texttt{iris\_inconnues} pour $k=3$. On initialise les centres de classes par la méthode heuristique décrite ci-dessus. 
}

\begin{quest}
  \QR	{Définir une fonction commentée \texttt{kmoy} qui prend comme arguments un entier $k$, un ensemble de données \texttt{E}, une distance \texttt{d}, et un nombre d'itérations $n$, et qui renvoie la classification des données de \texttt{E} obtenue par $n$ itérations de l'algorithme des $k$-moyennes. On pourra utiliser les fonctions \texttt{init\_centres},  \texttt{partition}, et \texttt{barycentre} définies précédemment. On pourra s'aider du code suivant.
    \inputminted[linenos=true,frame=single]{python}{\pathtoexos/Info/Scripts/ia_kmoy.py}}
  {
      \inputminted[linenos=true,frame=single]{python}{\pathtoexos/Info/Scripts/ia_kmoy_sol.py}
  }

  \QR {On rappelle ci-dessous à gauche la classification connue des variétés d'iris, que l'on a reporté de la \Fig{iris_graphe}. On donne ci-dessous à droite le classement obtenu par l'algorithme des $k$-moyennes pour $k=3$ sur la \Fig{iris_kmoy}. Commenter les résultats obtenus.}{

  }
\end{quest}

\Enonce{
\begin{fig}[iris_moy]
   \centering
   \begin{subfig}[][0.47]
     \includegraphics[height=0.7\columnwidth]{iris_graphe.png}
   \end{subfig}
   \begin{subfig}[][0.47]
     \includegraphics[height=0.7\columnwidth]{kmeans_iris_2.pdf}
   \end{subfig}
   \capt{
     (a) Classification connue des variétés d'iris.
     (b) Classement des variétés d'iris prédit par 
     l'algorithme des $k$-moyennes pour $k=3$. Le nombre d'itérations est choisi de sorte que les centres soient stationnaires, c'est-à-dire n'évoluent plus entre deux itérations. Les centres finaux sont représentés par des croix.
   }
 \end{fig}
}

 \begin{quest}
  \QR	{Pensez-vous que l'algorithme des $k$-moyennes est adapté pour classer un ensemble de données réparties comme sur la \Fig{kmeans_limit} ci-dessous ? Justifier.}{

  }
\end{quest}

\Enonce{
\begin{fig}[kmeans_limit]
  \centering
  \includegraphics[width=0.5\columnwidth]{kmeans_limit.png}
  \capt{Exemple de données non linéairement séparables. 
  % La classification par l'algorithme des $k$-moyennes est-elle adaptée ?
  }
\end{fig} 
}


 \subsection{Méthode du coude}

 \Enonce{
Nous avons imposé le nombre de classes à $k=3$ sachant que les données correspondent à l'une ou l'autre des trois variétés d'iris : setosa, versicolor et virginica. En présence de données brutes issues du terrain, et en particulier de plantes mal connues, le nombre de variétés de fleur n'est pas connu. On peut alors appliquer l'algorithme des $k$-moyennes pour différents nombres de classes $k$ et comparer les résultats. On définit les notions suivantes.

La \emph{variance} d'un ensemble de données $X$ de barycentre $b$ est la grandeur 
$$
V(X) = \sum_{x\in X} d(x,b)^2.
$$
L'\emph{inertie} d'une classification $p=(X_i)_{i=1,...,k}$ d'un ensemble de données $E$ en $k$ classes est la somme des variances des classes, soit
$$
I(p) = \sum_{i=1}^k V(X_i).
$$
 }

\begin{quest}
  \QR {\emph{Exemple}. Calculer les barycentres $b_1$ et $b_2$ et les variances $V_1$ et $V_2$ des classes $X_1$ et $X_2$ suivantes : $X_1 = \{(1,1),(3,3)\}$, $X_2=\{(5,4),(7,4)\}$. Calculer l'inertie $I$ de la classification $p=(X_1,X_2)$. Représenter ces données sur un graphe.}{

  }

  \QR	{Justifier que minimiser l'inertie $I(p)$ est un bon critère pour juger de la qualité d'un classement de données par une classification $p$. Comment évolue l'inertie $I(p)$ au cours des différentes itérations de l'algorithme des $k$-moyennes ? Commenter.}{

  }

  \QR	{On représente ci-dessous l'inertie de la classification obtenue par l'algorithme des $k$-moyennes en fonction du nombre de classes $k$. Pour déterminer la valeur optimale de $k$, il est d'usage de recourir à la \emph{méthode du coude} : on choisit la valeur de $k$ pour laquelle l'inertie diminue significativement avant cette valeur mais ne varie presque plus après cette valeur. Justifier ce critère. Quelle valeur de $k$ est optimale pour le classement des données d'iris d'après le graphe de la \Fig{elbow} ci-dessous ?}{
    
  }
\end{quest}

\Enonce{
\begin{fig}[elbow]
  \centering
  \includegraphics[width=0.5\columnwidth]{elbow.png}
  \capt{Inertie $I(p)$ de la classification $p$ obtenue par l'algorithme des $k$-moyennes sur les données d'iris pour différents nombres de classes $k$.}
\end{fig}
}

\end{exo}
%%%%%%%%%%%%%%%%%%%%%%%%%%%%%%%%%%%

%%%%%%%%%%%%%%%%%%%%%%%%%%%%%%%%%%%
\begin{exo}[PC][2][TD]{Apprentissage supervisé et classement d'articles de presse}
  Des informaticiens travaillant pour un média connu souhaitent élaborer un algorithme de classement des articles de presse pour le site internet du magazine. Les articles de presse se répartissent en 5 sujets : sport, politique, science, culture, économie. Pour que l'algorithme soit capable de classer un nouvel article de presse dans l'une de ces 5 classes, on procède par apprentissage supervisé et on recourt à l'algorithme des $k$ plus proches voisins. 

  \begin{fig}[label]
    \centering
    \includegraphics[width=0.4\columnwidth]{network.pdf}
    %\capt{}
  \end{fig}
  \parag{Classification} On dispose d'un ensemble d'apprentissage \texttt{articles} contenant les $N=280$ articles de presse déjà publiés. On souhaite créer une classification des articles de presse selon leur sujet. Pour cela, on modélise un article de presse par un point du plan $y=(e,n)$ où $0\leq e \leq 1$ représente la quantité d'équations dans le texte et $0\leq n \leq 1$ est la quantité de noms propres.
  
\begin{quest}
  \item	Rappeler le principe de l'apprentissage supervisé et d'un algorithme de classement.
  \item On donne le lien ci-dessous au fichier \texttt{article\_presse.py} contenant l'ensemble d'apprentissage  \texttt{articles} (liste de points) et la classification \texttt{sujets} (de type dictionnaire). Que vaut \texttt{sujets[(0.33, 0.31)]} ?
  \begin{center}
    \href{https://cahier-de-prepa.fr/pc-theo/download?id=682}{https://cahier-de-prepa.fr/pc-theo/download?id=682}
  \end{center}
  
  \item Définir en \texttt{python} une fonction distance \texttt{d} qui prend comme arguments deux points \texttt{y1} et \texttt{y2} du plan et qui renvoie la distance euclidienne entre ces deux points. 
\sol{
  \rm
      \inputminted[linenos=true,frame=single]{python}{\pathtoexos/Info/Scripts/ia_distance_euclid.py}
  }
\end{quest}

\parag{Phase d'apprentissage : algorithme des $k$ plus proches voisins} On souhaite implémenter puis appliquer l'algorithme des $k$ plus proches voisins pour classer le nouvel article $x=(0.5,0.6)$. On note de façon générique $E$ un ensemble d'apprentissage, $p$ une classification, $k$ entier le nombre de voisins considérés et $d$ la distance.
\begin{quest}
\item Rappeler le principe de l'algorithme des $k$ plus proches voisins. Que représentent les fonctions {\rm{tri}}$(E,x,d)$, {\rm{voisins}}$(E,x,d,k)$ et {\rm{classe}}$(E,x,d,k,p)$ ?
\item On souhaite implémenter la fonction {\rm{tri}}$(E,x,d)$. Rappeler les différents algorithmes de tri usuels. Compléter les lignes de code suivantes pour appliquer le tri bulle à l'ensemble d'apprentissage $E$. Vérifier cette fonction sur des exemples simples, par exemple pour \texttt{E=[(0.1,0.1),(0.9,0.9),(0.5,0.5)]} et \texttt{x=(0.4,0.6)}.
  \inputminted[linenos=true,frame=single]{python}{\pathtoexos/Info/Scripts/tri_xy.py} 
\sol{Tri bulle, par insertion, par sélection, rapide, fusion, comptage.
  \rm
      \inputminted[linenos=true,frame=single]{python}{\pathtoexos/Info/Scripts/tri_xy_sol.py}
}
\item Compléter le code de la fonction {\rm{voisins}}$(E,x,d,k)$ définie ci-dessous. 
  \inputminted[linenos=true,frame=single]{python}{\pathtoexos/Info/Scripts/voisins_xy.py}
\sol{\rm
    \inputminted[linenos=true,frame=single]{python}{\pathtoexos/Info/Scripts/voisins_xy_sol.py}
}
\item On détermine ensuite la classe dominante de {\rm{voisins}}$(E,x,d,k)$, c'est-à-dire la classe $i$ telle que l'intersection $E_j\, \cap $ {\rm{voisins}}$(E,x,d,k)$ contienne le maximum d'éléments pour $j=i$. Commenter chaque ligne du code suivant.
  \inputminted[linenos=true,frame=single]{python}{\pathtoexos/Info/Scripts/classe_xy.py}
Effectuer l'exécution pas à pas de l'appel \texttt{classemaj(E, x, d, k, p)} sur l'exemple simple \\ \texttt{E=[(0.2,0.2),(0.8,0.8),(0.5,0.5)]},
\texttt{p=\{(0.1,0.1):`culture',(0.9,0.9):`économie',}\\ 
\texttt{(0.5,0.5):`sport'\}}, \texttt{k=1} et \texttt{x=(0.4,0.6)}. 
\item En appelant la fonction \texttt{classemaj}, appliquer l'algorithme des $k$ plus proches voisins à l'ensemble d'apprentissage \texttt{articles} et à la classification \texttt{sujets} pour $x=(0.37,0.69)$ et $k=1,3,7$. Dans quel sujet classer l'article $x$ ? 

\end{quest}

\parag{Phase de test : matrice de confusion} On considère un ensemble test \texttt{articles\_test} contenant $N=20$ articles de presse à publier dont on connaît les sujets, répertoriés dans le dictionnaire \texttt{sujets}. 
% Voir lien ci-dessous.
% \begin{center}
%   \href{https://cahier-de-prepa.fr/pc-theo/download?id=682}{https://cahier-de-prepa.fr/pc-theo/download?id=682}
% \end{center}
On souhaite évaluer la qualité de l'algorithme des $k$ plus proches voisins en calculant la matrice de confusion $M$ associée à cet ensemble test. 
\begin{quest}
  \item	Rappeler la définition de la matrice de confusion $M$ associée à un ensemble test $T$ de classification $p$. 
  \sol{
    On rappelle que la matrice de confusion est une matrice $k\times k$ dont les lignes représentent les classes réelles et les colonnes les classes prédites. L'élément $M_{ij}$ de la matrice de confusion est le nombre d'articles de la classe $i$ classés dans la classe $j$.
  }
\item Créer un dictionnaire \texttt{numero} qui associe a chaque classe, de type \texttt{string} (chaîne de caractères), un entier de 0 à 4.
On donne l'implémentation \texttt{python} du calcul de la matrice de confusion ci-dessous. Compléter les lignes manquantes.
  \inputminted[linenos=true,frame=single]{python}{\pathtoexos/Info/Scripts/matrice_confus.py}
Appliquer cette fonction à l'ensemble test \texttt{articles\_test} puis inspecter la matrice de confusion.
\sol{
  \rm
      \inputminted[linenos=true,frame=single]{python}{\pathtoexos/Info/Scripts/matrice_confus_sol.py}
}
\item Implémenter les fonctions \texttt{precision} et \texttt{erreur}. Déterminer la précision P et le taux d'erreur TE.
\end{quest}

  % On considère un ensemble d'apprentissage $E$ contenant $N=100$ articles de presse. Ces articles sont répartis en deux classes : \{politique, sport\}. On souhaite apprendre à la machine à classer un article $x$ n'appartenant pas à $E$ grâce à l’algorithme des  $k$ plus proches voisins. On réalise donc un apprentissage supervisé. On utilise la distance euclidienne dans le plan, soit $d(x,y) = (y_1-x_1)^2 + (y_2-x_2)^2$.
  


\end{exo}
%%%%%%%%%%%%%%%%%%%%%%%%%%%%%%%%%%%

% %%%%%%%%%%%%%%%%%%%%%%%%%%%%%%%%%%%%%%%%%%%%%%%%%%%%%%%%%%%%%%%%
% \begin{exopython}[PC]{Partition et classification}{1}
% On considère un ensemble $E$ de mots de la langue française. On souhaite créer une classification de ces mots selon leur longueur : une classe pour les mots de $3$ lettres et moins ; une autre classe pour les mots de strictement plus de $3$ lettres.
% \begin{quest}
% \item Créer une liste \texttt{E} des mots \texttt{`or'}, \texttt{`tel'}, \texttt{`chat'}, \texttt{`appel'}, \texttt{`manque'}. De quel type python est l’objet \texttt{E} ? l’objet \texttt{`or'} ? L’objet \texttt{E} est-il mutable ? Que renvoie la commande \texttt{E[0]=`argent’} ?
% \item On souhaite créer la fonction \texttt{classement} qui classe les mots d’un ensemble \texttt{E} selon le critère ci-dessus. La fonction doit satisfaire les spécifications suivantes :
%   \inputminted[linenos=true,frame=single]{python}{\pathtoexos/Info/Scripts/spec_class.py}
% Rédiger sur papier un algorithme (avec la syntaxe du français) effectuant cette tâche.
% % Il pourra compléter les lignes ci-dessous

% %$\begin{array}{lll}
% %\Defin{classement($E$)} \\
% %& \Ini{$p =$ \{\}} \\
% %& \Ini{$n =$ \text{len}(E)}\\
% %& \For{i}{1}{n} \\
% %& & \Do{...} \\
% %\end{array}$
% %\item Démontrer mathématiquement que la classification $p$ que retourne l’algorithme forme bien une partition de $E$.

% \sol{
% La boucle \texttt{for} parcourt tous les éléments \texttt{mot} de $E$. À chaque étape de la boucle, la clé \texttt{mot} est nécessairement ajoutée. Donc la réunion des classes, c’est-à-dire des éléments de $p$, forme $E$.

% La condition \texttt{if} - \texttt{else} assure que la clé \texttt{mot} est soit dans la classe \texttt{`court'}, soit dans la classe \texttt{`long'}. Les classes sont donc disjointes. De toute façon, un dictionnaire supprime les doublons, donc un mot donné de $E$ ne peut appartenir qu’à une seule classe.
% }

% \item Écrire le programme python associé à cet algorithme, commenté et conforme aux spécifications. On nommera les classes par des chaînes de caractères, \texttt{`court'} si le mot fait moins de $3$ lettres, \texttt{`long'} sinon. Que renvoie ensuite le script suivant ?
%   \inputminted[linenos=true,frame=single]{python}{\pathtoexos/Info/Scripts/test_class.py}
% \sol{
%   \inputminted[linenos=true,frame=single]{python}{\pathtoexos/Info/Scripts/class_mots.py}
% }
% \end{quest}

% \end{exopython}
% %%%%%%%%%%%%%%%%%%%%%%%%%%%%%%%%%%%%%%%%%%%%%%%%%%%%%%%%%%%%%%%%


%%%%%%%%%%%%%%%%%%%%%%%%%%%%%%%%%%%%%%%%%%%%%%%%%%%%%%%%%%%%%%%%
\begin{exo}[PC][1][python]{Classification de points colorés du plan}
On considère un ensemble d'apprentissage $E$ contenant $N=100$ points du plan de coordonnées $y=(y_1,y_2)$. Ces données sont réparties en deux classes selon la couleur du point : '\{rouge, bleu\}.'. On souhaite apprendre à la machine à classer un point $x=(x_1,x_2)$ n'appartenant pas à $E$ grâce à l’algorithme des  $k$ plus proches voisins. On réalise donc un apprentissage supervisé. On utilise la distance euclidienne dans le plan, soit $d(x,y) = (y_1-x_1)^2 + (y_2-x_2)^2$. 
\begin{quest}
  \item On définit l'ensemble $E$ comme une liste de $N=100$ points du plan de position aléatoire et de coordonnées comprises entre $-1$ et $1$.
    \inputminted[linenos=true,frame=single]{python}{\pathtoexos/Info/Scripts/E_xy.py}
  À l'aide de la fonction \texttt{plot} de la bibliothèque \texttt{matplotlib.pyplot}, représenter ces points sur un graphe.
  \sol{  \inputminted[linenos=true,frame=single]{python}{\pathtoexos/Info/Scripts/E_xy_sol.py}}
  \item On partionne maintenant les points de $E$ selon leur couleur : rouge ou bleue.
  On colorise les points en rouge si leur distance $d$ à l'origine est supérieure à $\num{0.5}$, où la fonction distance $d$ est définie par
    \inputminted[linenos=true,frame=single]{python}{\pathtoexos/Info/Scripts/dist_xy.py}
  Définir la fonction \texttt{couleur} qui associe la bonne couleur à tout point $y$ du plan.
  Définir ensuite le dictionnaire \texttt{p} qui représente la classification des points $y$ de $E$ selon leur couleur.
  On définira les deux classes par les chaînes de caractères \texttt{`rouge'} et \texttt{`blue'}. Représenter les points sur un graphe avec leur bonne couleur.
  \sol{  \inputminted[linenos=true,frame=single]{python}{\pathtoexos/Info/Scripts/dist_xy_sol.py}}
  \item À partir des scripts de l'algorithme des k plus proches voisins vus en cours et téléchargeables sur Cahier de Prépa, déterminer la classe majoritaire du point $x=(\num{0.71};0)$ pour $k=1,3$ et $7$ voisins.
  \sol{  \inputminted[linenos=true,frame=single]{python}{\pathtoexos/Info/Scripts/xy_classe_sol.py}}
\end{quest}
\end{exo}
%%%%%%%%%%%%%%%%%%%%%%%%%%%%%%%%%%%%%%%%%%%%%%%%%%%%%%%%%%%%%%%%


%%%%%%%%%%%%%%%%%%%%%%%%%%%%%%%%%%%%%%%%%%%%%%%%%%%%%%%%%%%%%%%%
\begin{exo}[PC][1][python]{Qualité des prédictions}
On se donne une matrice de confusion $M$ représentée par une liste de listes d’entiers. Cette matrice caractérise la fiabilité d’un test médical : le test détermine si le patient est infecté ou sain. Le test a été effectué sur un nombre total de patients $N$.
\begin{quest}
\item Quelle est la taille de la matrice $M$ ? Comment calculer $N$ connaissant uniquement la matrice~$M$ ? Indiquer schématiquement ce que représente chaque élément de cette matrice.
\sol{
\begin{center}
\begin{tabular}{llll}
& & \textbf{Classe prédite} & \\
& & Infecté & Sain \\
\textbf{Classe réelle} & Infecté & $20$ (vrai positif) & $8$ (faux négatif)\\
 & Sain & $2$ (faux positif) & $70$ (vrai négatif) \\
\end{tabular}
\end{center}
}
\item Définir en \texttt{python} une fonction \texttt{sensibilite} qui prend comme seul argument  la matrice $M$ et qui calcule la sensibilité du test, c’est-à-dire le nombre de vrais positifs sur le nombre total de tests positifs. On précisera les spécifications de cette fonction et on la commentera. Tester cette fonction sur un exemple.
\item Définir en \texttt{python} une fonction \texttt{specificite}  qui prend comme seul argument  la matrice $M$ et qui calcule la spécificité du test, c’est-à-dire le nombre de vrais négatifs sur le nombre total de tests négatifs. On précisera les spécifications de cette fonction et on la commentera. Tester cette fonction sur un exemple.
\item Définir en \texttt{python} une fonction \texttt{prevalence} qui prend comme seul argument la matrice $M$ et qui calcule la prévalence du test, c’est-à-dire le nombre de patients infectées sur le nombre total de patients. On précisera les spécifications de cette fonction et on la commentera. Tester cette fonction sur un exemple.
\sol{
  \inputminted[linenos=true,frame=single]{python}{\pathtoexos/Info/Scripts/sens.py}
}
\end{quest}
\end{exo}
%%%%%%%%%%%%%%%%%%%%%%%%%%%%%%%%%%%%%%%%%%%%%%%%%%%%%%%%%%%%%%%%

%%%%%%%%%%%%%%%%%%%%%%%%%%%%%%%%%%%
\begin{exo}[PC][2][devoir]{Classement de fruits}
  On dispose d’un ensemble de données sur divers fruits caractérisés par leur masse et leur couleur. La masse est indiquée en grammes et représentée par un entier positif, tandis que la couleur est représentée par un flottant compris entre $0$ et $1$ qui indique si la couleur est plutôt \guill{chaude} ou \guill{froide}, suivant la correspondance suivante :
\begin{center}
\begin{tabular}{|c|c|}
  \hline
\textbf{Couleur} & \textbf{Valeur} \\
\hline
Rouge & $0.85$ - $1.00$ \\
Orange & $0.75$ - $0.85$ \\
Jaune & $0.65$ - $0.75$ \\
Vert & $0.45$ - $0.65$ \\
\hline
\end{tabular}
\end{center}
On souhaite classer ces de nouvelles données de fruits à l'aide de l'algorithme des $k$-moyennes, connaissant uniquement les données de masse et de couleur. On dispose pour cela d'un ensemble d'apprentissage, représentée par une liste \texttt{parametres} dont les premières entrées sont \texttt{parametres = [(123,0.81), (150,0.73), (100,0.65), ...]}, et d'une liste \texttt{fruits} de même longueur contenant les étiquettes de chaque fruit associées à \texttt{parametres}, et qui commence par \texttt{fruits = [`tomate', `banane', `kiwi', ...]}.  L'ensemble de ces données est représenté sur le graphique de la \Fig{fruits} ci-dessous. 

\begin{quest}
  \item	Le programmeur ne connaît pas encore l'ensemble des fruits représentés dans les données, et donc le nombre de classes à former. Reproduire, compléter et commenter la fonction \texttt{classes\_fruits} ci-dessous qui prend comme argument la liste \texttt{fruits} et qui renvoie l'ensemble des classes, i.e la liste de tous les fruits possibles. 
  
  \emph{Les pointillés ... peuvent représenter plusieurs lignes de code.}
  
\begin{minted}[linenos=true,frame=single]{python}
def classes_fruits(fruits):
    """
    Entrée : fruits (list) : liste de chaînes de caractères
    Sortie : classes (list) : liste de chaînes de caractères
    """
    classes = []
    ...
    return classes
\end{minted}  



\sol{
\inputminted[linenos=true,frame=single]{python}{\pathtoexos/Info/Scripts/fruits_0.py}
}

  \item Reproduire, compléter et commenter la fonction \texttt{classification} ci-dessous qui prend comme arguments les listes \texttt{fruits} et \texttt{parametres}, et qui renvoie un dictionnaire contenant le classification de l'ensemble d'apprentissage, i.e. tel que chaque entrée est de la forme \texttt{(masse, couleur) : fruit}.  
  
    \emph{Les pointillés ... peuvent représenter plusieurs lignes de code.}


    \begin{minted}[linenos=true,frame=single]{python}
def classification(fruits, parametres):
    """
    Entrée : fruits (list) : liste de chaînes de caractères
            parametres (list) : liste de tuples (masse, couleur)
    Sortie : p (dict) : dictionnaire contenant la classification
    """
    p = {}
    ...
    return p
  \end{minted}

\sol{
    \inputminted[linenos=true,frame=single]{python}{\pathtoexos/Info/Scripts/fruits_1.py}
}
\end{quest}


  % \begin{fig}
  % \centering
  %     \begin{subfig}[label][0.49]
  %         \includegraphics[width=0.47\textwidth]{fruits_-1.pdf}
  %     \end{subfig}
  % \begin{subfig}[label][0.49]
  %     \includegraphics[width=0.47\textwidth]{fruits_1.pdf}
  % \end{subfig}
  % \begin{subfig}[label][0.49]
  %     \includegraphics[width=0.47\textwidth]{fruits_8.pdf}
  % \end{subfig}
  % \begin{subfig}[label][0.49]
  %     \includegraphics[width=0.47\textwidth]{fruits_10.pdf}
  % \end{subfig}
  % \capt{(a) Jeu de données initial. L’abscisse correspond à la masse d’un fruit, l’ordonnée à la quantité de vert dans sa couleur (en RGB codé sur $8$ bits, soit $256$ valeurs). Ces données représentent trois fruits :  des tomates, des bananes, et des kiwis. (b) Partition initiale calculée à partir de centres aléatoires. Les losanges représentent les centres. (c) Partition après $8$ ré-évaluations des centres selon l’algorithme des $k$ moyennes. (d) Partition après $10$ ré-évaluations.}
  % \end{fig}
  
  
  Pour la distance, nous prenons une distance euclidienne dans l’espace (masse, couleur). Comme ces deux grandeurs sont de dimensions différentes, on peut pondérer ces deux grandeurs par des coefficients différents, de façon à donner un poids équitable à chaque paramètre. La couleur est déjà normalisée entre $0$ et $1$, tandis que la masse peut varier entre $0$ et $300$ grammes typiquement. Nous normaliserons donc la masse par l'unité de masse maximale fixée conventionnellement à $300$ grammes.
  \begin{quest}
    \item	Reproduire, compléter et commenter la fonction \texttt{distance} ci-dessous qui prend comme arguments deux points du plan  de coordonnées $(m_1,c_1)$ et $(m_2,c_2)$, et qui renvoie la distance euclidienne pondérée entre ces deux points.
     
     \begin{minted}[linenos=true,frame=single]{python}
def distance(p1, p2):
    """
    Entrée : p1 (tuple) : un point du plan de coordonnées (masse, couleur)
            p2 (tuple) : un point du plan de coordonnées (masse, couleur)
    Sortie : distance euclidienne pondérée entre p1 et p2
    """
    m1, c1 = p1
    m2, c2 = p2
    return ...
    \end{minted}


\sol{
      \inputminted[linenos=true,frame=single]{python}{\pathtoexos/Info/Scripts/fruits_distance_xy_sol.py}
}
  \end{quest}

  Nous appliquons maintenant l'algorithme des $k$ plus proches voisins.
  \begin{quest}
    \item	Reproduire, compléter et commenter la fonction \texttt{kvoisins} ci-dessous, qui prend comme arguments un point du plan de coordonnées $(m,c)$, l'ensemble d'apprentissage \texttt{parametres}, la classification \texttt{p}, la distance \texttt{distance}, et un entier $k$, et qui renvoie la classe majoritaire parmi les $k$ plus proches voisins de ce point. On précisera les spécifications de cette fonction et on la commentera.

    \emph{Les pointillés ... peuvent représenter plusieurs lignes de code.}


    \begin{minted}[linenos=true,frame=single]{python}
def kvoisins(point, parametres, p, distance, k):
    """
    Entrée : point (tuple) : point du plan (masse, couleur)
        parametres (list) : liste de tuples
        p (dict) : dictionnaire contenant la classification
        distance (func) : fonction distance entre deux points du plan
        k (int) : nombre de plus proches voisins
    Sortie : classe (string) : classe majoritaire parmi les k plus proches voisins
    """
    tri = sorted(parametres, key=lambda p: distance(point, p))
    voisins = tri[:k]
    classes_count = {}
    for voisin in voisins:
        ...
    classe_majoritaire = max(classes_count, key=lambda c: classes_count[c])
    return classe_majoritaire
    \end{minted}


\sol{
\inputminted[linenos=true,frame=single]{python}{\pathtoexos/Info/Scripts/fruits_2.py}
}

    \item On représente l'ensemble d'apprentissage \texttt{parametres} sur la \Fig{fruits} ci-dessous, pour un exemple de trois fruits : des tomates, des bananes, et des kiwis. On applique la fonction \texttt{kvoisins} à un fruit de coordonnées normalisées $x=(0.77,0.78)$. En s'aidant d'un compas ou d'une règle, déterminer graphiquement la classe majoritaire de cette nouvelle donnée pour $k=1,3$ et $7$ plus proches voisins, en expliquant votre construction. Commenter les résultats obtenus. 
    
    \emph{Pensez à rendre l'énoncé avec le graphique complété avec votre copie !}

    \sol{
    Pour $k=1$, le voisin le plus proche de $x$ est un point de classe \guill{Tomate}, donc la classe majoritaire est \guill{Tomate}.


    Pour $k=3$, les trois voisins les plus proches de $x$ sont deux points de classe \guill{Banane} et un point de classe \guill{Tomate}, donc la classe majoritaire est \guill{Banane}.


    Pour $k=7$, les sept voisins les plus proches de $x$ sont quatre points de classe \guill{Tomate} et trois points de classe \guill{Banane}, donc la classe majoritaire est \guill{Tomate}.

    On remarque que le résultat de l'algorithme des $k$ plus proches voisins dépend du choix de $k$ !
    }
  \end{quest}
  \begin{fig}[fruits]
 \centering
 \includegraphics[width=\textwidth]{fruits_10.pdf}
 \capt{Ensemble d'apprentissage de fruits représenté dans le plan (masse, couleur). }
 \end{fig}

 On souhaite améliorer la prédiction de l'algorithme des $k$ plus proches voisins en pondérant la classe majoritaire par le nombre de voisins appartenant à chaque classe.le classes des voisins suivant la distance au nouveau point $x$. Notons $V_i(x,k)$ le sous-ensemble des voisins de $x$ de classe $i$ parmi les $k$ plus proches voisins de $x$. On propose d'attribuer à chaque classe  un poids égal  à
  \begin{equation}
    \sigma_i = \sum_{v \in V_i(x,k)} d(x, v)^{- 2}.
  \end{equation}
  
 \begin{quest}
  \item	
  On donne la liste triée des $7$ premiers voisins de la donnée $x=(0.77,0.78)$ ci-dessous, avec leur classe et leurs coordonnées normalisées. Calculer les poids $\sigma_i$ de chaque classe \guill{Tomate} ou \guill{Banane} pour les $k=1,3$ et $7$ plus proches voisins de $x$.

  \sol{
    % distance d and poids sigma for each point v in V_i(x,k)
    \begin{center}
      \begin{tabular}{|c|c|c|c|c|c|c|c|}
        \hline
        Voisin & 1 & 2 & 3 & 4 & 5 & 6 & 7 \\
        \hline
        Classe & Tomate & Banane & Banane & Tomate & Tomate & Tomate & Banane \\
        \hline
        Coordonnées & $(0.76,0.80)$ &  $(0.79,0.80)$ & $(0.78,0.75)$ & $(0.75,0.84)$ & $(0.76,0.73)$ & $(0.74,0.84)$ & $(0.82,0.80)$\\
        \hline
        $d(v, x)^{-2}$ & $2000$ & $1250$ & $1000$ & $250$ & $385$ & $222$ & $345$ \\
        \hline
      \end{tabular}
    \end{center}

    On calcule le poids des classes pour chaque valeur de $k$.
    \begin{center}
      \begin{tabular}{|c|c|c|}
        \hline
        $k$ & $\sigma_{\text{Tomate}}$ & $\sigma_{\text{Banane}}$ \\
        \hline
        1 & $2000$ & $0$ \\
        \hline
        3 & $2000$ & $2250$ \\
        \hline
        7 & $2857$ & $2595$ \\
        \hline
      \end{tabular}
    \end{center}
  }
  \begin{center}\small
  \begin{tabular}{|c|c|c|c|c|c|c|c|}
    \hline
    Voisin & 1 & 2 & 3 & 4 & 5 & 6 & 7\\
    \hline
    Classe & Tomate & Banane & Banane & Tomate & Tomate & Tomate & Banane \\
    \hline
    Coordonnées & $(0.76,0.80)$ &  $(0.79,0.80)$ & $(0.78,0.75)$ & $(0.75,0.84)$ & $(0.76,0.73)$ & $(0.74,0.84)$ & $(0.82,0.80)$\\
    \hline
  \end{tabular}
  \end{center}
  \item On choisit d'attribuer à la nouvelle donnée $x$ la classe dont le  poids est le plus élevé. Quelle classe est attribuée à $x$ pour $k=1,3$ et $7$ voisins ? Commenter les résultats obtenus.
  \sol{
  On trouve comme résultats :
\begin{center}
  \begin{tabular}{|c|c|}
    \hline
    $k$ & Classe attribuée à $x$ \\
    \hline
    1 & Tomate \\
    \hline
    3 & Banane \\ 
    \hline
    7 & Tomate \\
    \hline
  \end{tabular}
\end{center}
On voit que pour $k=3$, les tomates et les bananes ont des poids assez proches, mais toujours au profit des bananes. Malgré l'amélioration de la prédiction grâce à la pondération, le résultat de l'algorithme des $k$ plus proches voisins dépend toujours du choix de $k$. La nouvelle donnée $x$ est située à la frontière entre les classes \guill{Tomate} et \guill{Banane}, ce qui rend sa classification plus délicate.
  }
 \end{quest}
  
  
\end{exo}
%%%%%%%%%%%%%%%%%%%%%%%%%%%%%%%%%%%